\begingroup
\setlength{\LTcapwidth}{\textwidth}
\begin{longtable}{p{0.37\textwidth}p{0.49\textwidth}}
\caption{Konfigurasi utama implementasi sistem.}
\label{tab:tab5.3_konfigurasi_utama_implementasi} \\
\toprule
\textbf{Komponen} & \textbf{Konfigurasi} \\
\midrule
\endfirsthead

\multicolumn{2}{@{}l}{\small\emph{Tabel \thetable{} (lanjutan)}} \\
\toprule
\textbf{Komponen} & \textbf{Konfigurasi} \\
\midrule
\endhead

\bottomrule
\endlastfoot

Model produksi & HistGradientBoostingRegressor \\

Random seed & 42 \\

Hyperparameter utama & Default scikit-learn, selain \texttt{random\_state} \\

Normalisasi & StandardScaler, fit pada data pelatihan \\

Basis pengetahuan & Koleksi \texttt{diabetes\_kb} \\

Chunk size & 900 karakter \\

Chunk overlap & 120 karakter \\

Batas kalimat & Diaktifkan pada proses chunking produksi \\

Jumlah potongan terindeks & 2.233 potongan \\

Retriever produksi & Okapi BM25 \\

Jumlah evidence & \texttt{top\_k = 5} \\

Temperature generasi & 0,2 \\
\end{longtable}
\endgroup
